% 原生矢量对象；文字与对象在同一TeX画布排版。%
% 可选绘图样式用于被点名对象的透明表面；默认空值保留既有填充。
\providecommand{\trustVectorRecord}[4][]{%
 \begin{scope}[shift={(#2,#3)}]%
  \path[draw=FigureStroke,fill=white,rounded corners=.5mm,#1] (-3.8,-5)--(3.8,-5)--(3.8,2.5)--(1.3,5)--(-3.8,5)--cycle;%
  \draw[FigureStroke] (1.3,5)--(1.3,2.5)--(3.8,2.5);%
  \foreach \r in {-2.8,-.8,1.2}{\draw[#4,line width=.9pt] (-2.3,\r)--(2.2,\r);}%
 \end{scope}}%
\providecommand{\trustVectorStore}[4][]{%
 \begin{scope}[shift={(#2,#3)}]%
  \path[draw=FigureStroke,fill=#4!12,#1] (-5,5)--(-5,-5) arc[start angle=180,end angle=360,x radius=5,y radius=1.8]--(5,5);%
  \draw[FigureStroke] (-5,0) arc[start angle=180,end angle=360,x radius=5,y radius=1.8];%
  \draw[FigureStroke] (-5,-2.5) arc[start angle=180,end angle=360,x radius=5,y radius=1.8];%
  \path[draw=FigureStroke,fill=#4!20,#1] (0,5) ellipse[x radius=5,y radius=1.8];%
 \end{scope}}%
\providecommand{\trustVectorMail}[5][]{%
 \begin{scope}[shift={(#2,#3)}]%
  \path[draw=FigureStroke,fill=#4!10,rounded corners=1mm,#1] (-6,-4) rectangle(6,4);%
  \draw[#4,line width=1pt] (-5.5,3.5)--(0,-.5)--(5.5,3.5);%
  \node[font=\figurefont\fontsize{7.5}{9}\selectfont,fill=white,inner sep=.6pt,#1] at (4.7,-3.5) {#5};%
 \end{scope}}%
\providecommand{\trustVectorNetwork}[3]{%
 \begin{scope}[shift={(#1,#2)}]%
  \foreach \a in {-5,0,5}{\foreach \b in {-4,4}{\draw[FigureMuted,line width=.55pt] (-8,\a)--(0,\b);}}%
  \foreach \a in {-4,4}{\foreach \b in {-3,3}{\draw[FigureMuted,line width=.55pt] (0,\a)--(8,\b);}}%
  % 第三参数保留旧调用接口；神经元按层实色填充蓝、红、黄。
  \foreach \xx/\yy/\col in {-8/-5/FigureBlue,-8/0/FigureBlue,-8/5/FigureBlue,0/-4/FigureRed,0/4/FigureRed,8/-3/FigureYellow,8/3/FigureYellow}{\path[draw=FigureStroke,fill=\col,line width=.75pt] (\xx,\yy) circle(1.7);}%
 \end{scope}}%
\providecommand{\trustVectorKey}[3]{%
 \begin{scope}[shift={(#1,#2)}]%
  \draw[#3,line width=1.2pt] (-3,0) circle(2.4);%
  \draw[#3,line width=1.2pt] (-.6,0)--(6,0) (3,0)--(3,-2) (5,0)--(5,-2);%
 \end{scope}}%
\providecommand{\trustVectorLock}[4][]{%
 \begin{scope}[shift={(#2,#3)}]%
  \draw[#4,line width=.85pt,rounded corners=1.5mm] (-2.2,1)--(-2.2,4)--(2.2,4)--(2.2,1);%
  \path[draw=FigureStroke,fill=white,rounded corners=.6mm,#1] (-3,-3) rectangle (3,1.5);%
  \fill[#4] (0,-.2) circle(.7); \draw[#4] (0,-.2)--(0,-1.8);%
 \end{scope}}%
\providecommand{\trustVectorPackage}[3]{%
 \begin{scope}[shift={(#1,#2)}]%
  \path[draw=FigureStroke,fill=#3!12] (-5,1)--(0,4)--(5,1)--(5,-5)--(0,-8)--(-5,-5)--cycle;%
  \draw[FigureStroke] (-5,1)--(0,-2)--(5,1) (0,-2)--(0,-8);%
  \draw[#3,line width=1.2pt] (-2.5,2.5)--(2.5,-.5);%
 \end{scope}}%
\providecommand{\trustVectorStop}[3]{%
 \begin{scope}[shift={(#1,#2)}]%
  \fill[white] (-1,-4) rectangle (1,4);%
  \draw[#3,line width=1.8pt] (0,-3.4)--(0,3.4);%
 \end{scope}}%
\providecommand{\trustVectorMagnifier}[3]{%
 \begin{scope}[shift={(#1,#2)}]%
  \draw[#3,line width=1.1pt] (0,0) circle(3);%
  \draw[#3,line width=1.4pt] (2.2,-2.2)--(5,-5);%
 \end{scope}}%
